\subsection{Qué significa exactamente «información de $n$ saltos»}
\label{sec:informacion}

Afirmar que un mecanismo es distribuido exige decir qué puede \emph{calcular}
cada agente, no solo con quién habla. Se adopta para ello la noción de grafo de
proximidad espacialmente distribuido \cite{bullo2009distributed}.

\begin{definicion}[Grafo espacialmente distribuido]
\label{def:espacialmente}
Un grafo de proximidad $\Gcal_1$ es \emph{espacialmente distribuido} sobre
$\Gcal_2$ si, para toda configuración $P$ y todo nodo $p$,
\[
  \Ncal_{\Gcal_1,p}(P)=\Ncal_{\Gcal_1,p}\bigl(\Ncal_{\Gcal_2,p}(P)\bigr),
\]
es decir, si un nodo que conoce la posición de sus vecinos en $\Gcal_2$ puede
calcular sus vecinos en $\Gcal_1$ sin más información.
\end{definicion}

El grafo de comunicación es el de disco $\Gcal_{\mathrm{disk}}(R^{\mathrm{com}})$.
La pregunta relevante es si el grafo de interacción del juego —el que define qué
pares comparten barrera y qué desviaciones son conectadas— puede construirse a
partir de lo que la radio entrega.

\begin{proposicion}[Condición de implementabilidad]
\label{pr:implementable}
Un mecanismo cuyo grafo de interacción $\Gcal_I$ sea espacialmente distribuido
sobre $\Gcal_{\mathrm{disk}}(R^{\mathrm{com}})$ es implementable con
información de un salto. Los grafos de vecindad relativa y de Gabriel
intersecados con el disco, y el grafo de Delaunay $r$-limitado
$\Gcal_{\mathrm{LD}}(R^{\mathrm{com}})$, cumplen esa
condición \cite[Prop.~2.9]{bullo2009distributed}; el grafo de Delaunay completo
intersecado con el disco, no.
\end{proposicion}

\begin{observacion}[Consecuencia de diseño]
\label{obs:eleccion-grafo}
La elección natural para $\Gcal_I$ no es el propio grafo de disco sino
$\Gcal_{\mathrm{LD}}(R^{\mathrm{com}})$: es espacialmente distribuido sobre él,
por tanto implementable, y estrictamente más disperso, lo que reduce el número
de pares que intercambian precio y el número de pares que deben
acordar un precio en \cref{th:precio-seguridad}. Un mecanismo que use un grafo
de interacción no espacialmente distribuido sobre el de comunicación no es
distribuido, aunque su descripción lo sugiera.
\end{observacion}

La elección puede afinarse más, porque los grafos de proximidad se ordenan y ese
orden tiene consecuencias operativas exactas.

\begin{teorema}[Orden, conectividad y localidad]
\label{th:jerarquia-grafos}
Sean $\Gcal_{\mathrm{EMST}}$, $\Gcal_{\mathrm{RN}}$, $\Gcal_{\mathrm{G}}$ y
$\Gcal_{\mathrm{D}}$ el árbol generador mínimo euclídeo, el grafo de vecindad
relativa, el de Gabriel y el de Delaunay \cite{bullo2009distributed}. Entonces:
\begin{enumerate}[label=(\roman*),leftmargin=2.2em,topsep=0.2em,itemsep=0.15em]
  \item $\Gcal_{\mathrm{EMST}}\subset\Gcal_{\mathrm{RN}}\subset
  \Gcal_{\mathrm{G}}\subset\Gcal_{\mathrm{D}}$;
  \item $\Gcal_{\mathrm{RN}}\cap\Gcal_{\mathrm{disco}}(R^{\mathrm{com}})$,
  $\Gcal_{\mathrm{G}}\cap\Gcal_{\mathrm{disco}}(R^{\mathrm{com}})$ y
  $\Gcal_{\mathrm{LD}}(R^{\mathrm{com}})$ tienen \emph{exactamente} las mismas
  componentes conexas que el grafo de disco;
  \item los tres son espacialmente distribuidos sobre el grafo de disco, es
  decir, cada AMR calcula su vecindad con la información de sus vecinos de
  disco;
  \item en cambio $\Gcal_{\mathrm{D}}\cap\Gcal_{\mathrm{disco}}$ es subgrafo del
  de disco pero \emph{no} es espacialmente distribuido sobre él.
\end{enumerate}
\end{teorema}

\begin{observacion}[Qué compra esto y a qué precio]
\label{obs:esparsificar}
La combinación de (ii) y (iii) es lo que interesa: se puede adelgazar el grafo de
comunicación sin perder \emph{ninguna} propiedad de conectividad, y cada AMR
puede decidir por sí mismo qué aristas conserva. Sobre veintidós AMR con
posiciones aleatorias, el grafo de Gabriel cortado al disco reduce las aristas
entre un $40\%$ y un $70\%$ según el radio, y el número de componentes conexas
coincide exactamente en los tres casos ensayados; la comprobación de (iii) no
encuentra ninguna discrepancia entre el cálculo local y el global.

El ahorro se compone con el de \cref{obs:etm-ahorro} porque actúa sobre un factor
distinto: aquel reduce \emph{cuántas veces} se habla, este reduce \emph{con
cuántos}. El precio es el diámetro: adelgazar alarga los caminos, y por
\cref{th:conmutada} la tasa de mezcla empeora. Se cambian aristas por rondas, y
qué convenga depende de si el cuello de botella es el ancho de banda o la
latencia.

El apartado (iv) es la advertencia: la restricción de Delaunay al disco parece la
elección natural —es la dual del diagrama de Voronoi que aparece en
\cref{th:potencia}— y sin embargo no es localmente calculable. Usarla exigiría
información que el grafo de comunicación no entrega.
\end{observacion}

\begin{observacion}[Cuatro grafos que no son el mismo]
\label{obs:cuatro-grafos}
Una precisión sobre el alcance de \cref{obs:eleccion-grafo}.
Que un grafo sea espacialmente distribuido justifica \emph{conocer a los
vecinos} con información local; no justifica calcular con ella todo el pago
—que puede depender de estados, mapas o restricciones que no viven en esa
bola— ni, sobre todo, \emph{omitir} restricciones de colisión. Con huellas
heterogéneas, la ausencia de una arista de radio entre dos AMR no excluye la
interferencia entre las cargas grandes que transportan. Omitir un par exige
demostrar que es seguro, o que otra restricción implica su seguridad; la
dispersión del grafo de radio no es un sustituto de la geometría
\cite{mayorga2026m1}. El grafo de comunicación, el de contactos, el de
conflictos y el de revisiones son objetos distintos, y el adelgazamiento de
\cref{obs:esparsificar} actúa sobre el primero solamente.
\end{observacion}

\figGrafos

Con $n>1$ saltos la condición se relaja: basta que $\Gcal_I$ sea espacialmente
distribuido sobre la potencia $n$-ésima del grafo de disco, lo que amplía la
clase admisible a cambio de $n$ rondas de intercambio y de una latencia
proporcional.
